{\color{red}
\section{Open questions for the author (NOTE list)}
\noindent[\textbf{New section.} Points that the fix notes do not determine uniquely, or that are not settled; in each case the text of the file was left as the author has it (or as the fix note gives it, where the fix note is unambiguous), with a red \textbf{NOTE} at the place.]
\begin{enumerate}
\item \textbf{DGLAP subtraction term} (DGLAP section). fix\_poles\_dglap Step 11 writes $-\frac{2N_{c}}{\xi}\mathcal{F}(0)$; this equals the author's subtraction $-2N_{c}\frac{1-\xi}{\xi}\mathcal{F}(0)$ \emph{plus} the remnant of Group II Row IV, while Step 11 also keeps the remnant in Row IV. The author's convention (remnant in Row IV) is kept. Which writing is wanted for the paper?
\item \textbf{Fragmentation-function scheme} (DGLAP section). The author's $\log\frac{(y'-y)^{2}\mu^{2}}{4\xi^{2}}$ is kept; it corresponds to $\overline{\rm MS}$ number-density FFs renormalised with the $d$-dimensional Jacobian $\xi^{-d}$. The CXY form differs by $\Delta$, Eq.~\eqref{eq:fx-Delta}. The author should confirm the FF definition used; no independent comparison with a standard NLO collinear $pp\to gX$ calculation was made.
\item \textbf{Coordinate or local form of DGLAP and RC} (DGLAP and RC sections). The coordinate forms are each logarithmically IR divergent at large parent position; the local form (Eq.~\eqref{eq:fx-loc}) is not. Choice for the paper: the author's (pending fix\_tail).
\item \textbf{Group II Row IV large-distance tail} (Group II Rows I, IV; Group III Row III). Compensation by Group II Row I (second part) and Group III Row III is established at the amplitude level only; the term-by-term check is pending (\texttt{notes/fix\_tail.tex}).
\item \textbf{Prescription of the energy denominators of $\mf{B}_{3}$} (Group I Row III; block $\mathcal{S}$). Principal value assumed (as in the setup); a $\pm i\epsilon$ prescription would add on-shell terms to the separate equations; not evaluated.
\item \textbf{Order of limits} (Finite terms, intro). The remaining pole $-\frac{\ell}{2}\frac{2}{\epsilon}$, $\ell=\log(1-\frac{k^{+}}{\vee})$, vanishes only if $\vee\to\infty$ is taken before $\epsilon\to0$; this is assumed and should be stated in the paper.
\item \textbf{Cubic contact terms} (singlet section). The cubic rows leave quadratic contact terms after the projectile average; whether they cancel among the rows, or against the contact terms of the three-$\rho$ evolution terms, was not checked.
\item \textbf{Integrability at the contact configurations} (singlet section): checked only for two kernels.
\item \textbf{Assumption (A2)} (singlet section): no C-odd ($d^{abc}$) or chiral ($f^{abc}g_{f}$) three-point function of the projectile; otherwise the symmetrised cubic parts, odd in $k$, are added.
\item \textbf{Kinematically constrained reading} of the soft block (fix\_soft Step 7): equivalent, not used; the author may prefer it for the paper.
\item \textbf{Not done in fix\_soft}: complete numerical integration with Wilson lines; re-check of the regular terms (taken from \texttt{Finite\_terms\_final}); the non-singlet $\log\frac{k^{+}}{\Lambda}\int\frac{d^{2}z}{z^{2}}C_{A}$ of the JIMWLK formula.
\end{enumerate}}\par
